\documentclass[11pt,leqno]{article}
\usepackage{fontspec}
\usepackage[french]{babel}
\usepackage{amsmath}
\usepackage{amssymb}
\usepackage{amsthm}
\usepackage{mathrsfs}
\usepackage{tikz-cd}
\usepackage{enumitem}

\frenchsetup{AutoSpacePunctuation=false}

\newcommand{\printedpage}[1]{\marginpar{\footnotesize p.~#1}}
\newcommand{\msp}[1]{\marginpar{\footnotesize [#1]}}
\newcommand{\DD}{\mathbb{D}}
\newcommand{\cA}{\mathcal{A}}
\newcommand{\Dia}{\mathsf{Dia}}
\newcommand{\Ob}{\mathsf{Ob}}
\newcommand{\Hom}{\mathsf{Hom}}
\newcommand{\hot}{\mathsf{hot}}
\newcommand{\Hot}{\mathsf{Hot}}
\newcommand{\HOT}{\mathsf{HOT}}
\newcommand{\isoto}{\overset{\sim}{\to}}
\newcommand{\isolongto}{\overset{\sim}{\longrightarrow}}

\begin{document}

\printedpage{34}
% head: A. Grothendieck : Dérivateurs, Ch. I, Édition des Universités de Montpellier II et Paris VII
% foot: Version 8 juillet 2001
\noindent est à fibres contractiles.

Autre situation standard où le changement de base marche (mais où les foncteurs qui interviennent n'ont aucune raison d'être propres ou lisses).

Soient
\[
f : I \longrightarrow I' \quad , \qquad g : J \longrightarrow J'
\]
dans $\Dia$, d'où le diagramme (comparer page 23)
\[
\begin{tikzcd}
X' = I \times J \arrow[r, "G'"] \arrow[d, "F'"'] & I \times J' = X \arrow[d, "F"] \\
Y' = I' \times J \arrow[r, "G"] & I' \times J' = Y
\end{tikzcd}
\quad .
\]
Ce diagramme est tel que pour tout dérivateur à gauche

\noindent\msp{page 60}

\noindent on a le changement de base
\[
G^*F_* \simeq F'_*{G'}^* \quad ,
\]
la raison ici étant que les foncteurs induits
\[
X'/y' \longrightarrow X/y
\]
sont des hot-équivalences \underline{\underline{co}}strictes, \textit{i.e.} que les catégories $K$ (page 49) sont contractiles. En effet, si $x_0 = (i_0, j'_0) \in \Ob X$, $y'_0 = (i'_0, j_0) \in \Ob Y'$, $u_0 = (u'_0, v'_0) : F(x_0) \to G(y'_0)$, où $u'_0 : f(i_0) \to i'_0$, $v'_0 : j'_0 \to g(j_0)$ sont donnés, on vérifie que la catégorie $K$ cherchée est isomorphe simplement à $(i_0, u'_0)\backslash(I/i'_0) \times (j_0, v'_0)\backslash(J/j'_0)$ qui est contractile ayant un objet initial \texttt{[plutôt $(i_0, u'_0)\backslash(I/i'_0) \times (j'_0\backslash J)/(j_0, v'_0)$ qui est contractile étant produit d'une catégorie ayant un objet initial et d'une catégorie ayant un objet final]}. Ouf!

Je m'aperçois d'un très gros canular. Je croyais me rappeler que pour $f : X \to Y$ qui est une hot-équivalence stricte, on avait pour tout $F \in \DD(Y)$, $F \isoto f_*f^*F$ (par exemple pour $\DD = \mathsf{Der}(\mathit{Ab})$, ou $\DD = \Hot$), \textit{i.e.} $f^*$ pleinement fidèle, d'où aussi que $\Gamma_{\DD}(Y, F) \to \Gamma_{\DD}(X, f^*F)$ est \texttt{[un]} isomorphisme. Mais ce n'est vrai que si $F$ est un coefficient \textit{constant} sur $Y$, \textit{i.e.} de la forme $q^*(M)$ ($q : Y \to e$, $M \in \DD(e) = \cA_{\DD}$). C'est faux, pour $F$ général, déjà dans le cas de l'immersion $f : \{\varphi\} \to Y$, où $\varphi$ est un objet initial. Alors, $f$ est toujours une hot-équivalence stricte, mais l'homomorphisme $\Gamma(Y, F) \to \Gamma(X, f^*F)$ n'est en général pas un isomorphisme.

\clearpage
\printedpage{35}
% head: A. Grothendieck : Dérivateurs, Ch. I, Édition des Universités de Montpellier II et Paris VII
% foot: Version 8 juillet 2001
\noindent\msp{page 61}

\subsection*{Notions de $\mathbb{D}$-équivalence et de $\mathbb{D}$-contractibilité.}

On va donc reprendre les questions d'acyclicité et de changement de base.

\noindent\textbf{Définition 1.} Soit $\DD$ un prédérivateur, et soit
\[
\begin{tikzcd}
X \arrow[rr, "f"] \arrow[dr, "p"'] & & Y \arrow[dl, "q"] \\
& e &
\end{tikzcd}
\]
dans $\Dia$. On dit que $f$ est une $\DD$-\textit{équivalence} si le foncteur
\[
F \longmapsto f^*F \quad , \qquad \DD(Y) \longrightarrow \DD(X)
\]
induit un foncteur pleinement fidèle sur la sous-catégorie de $\DD(Y)$ formée des $\DD$-coefficients \textit{constants}, \textit{i.e.} si pour tous $M, N \in \Ob \cA = \DD(e)$ (“coefficients absolus”), on a
\[
\Hom(q^*M, q^*N) \isolongto \Hom(p^*M, p^*N) \quad . \tag*{(1)}
\]
Supposons que $p_*$ et $q_*$ existent. On posera, pour un faisceau $F$ sur $X$
\[
H^*_{\DD}(X, F) \overset{\text{déf}}{=} p_*(F) \qquad \text{(et de même } \ H^{\DD}_*(X, F) = p_!(F) \ \text{).}
\]
N.B. Et si $F = p^*M$, $M \in \cA$, on note simplement $H^*_{\DD}(X, M)$, $H^{\DD}_*(X, M)$.

\noindent L'homomorphisme (1) s'interprète par adjonction comme
\[
\begin{tikzcd}[row sep=small]
\Hom(M, q_*q^*N) \arrow[r] \arrow[d, equal] & \Hom(M, p_*p^*N) \arrow[d, equal] \\
\Hom(M, H^*_{\DD}(Y, q^*N)) \arrow[r] & \Hom(M, H^*_{\DD}(X, p^*N))
\end{tikzcd}
\]
déduit de l'homomorphisme
\[
q_*q^*N \longrightarrow p_*p^*N \quad , \qquad \textit{i.e.} \quad H^*_{\DD}(Y, N) \longrightarrow H^*_{\DD}(X, N) \quad . \tag*{(2)}
\]
Donc dire que (1) est un isomorphisme pour tous $M, N$ signifie que (2) est un isomorphisme pour tout $N$. De même, si $p_!$, $q_!$ existent cela équivaut à dire que
\[
p_!p^*N \longrightarrow q_!q^*N \quad , \qquad \textit{i.e.} \quad H^{\DD}_*(X, N) \longrightarrow H^{\DD}_*(Y, N) \tag*{$(2')$}
\]
est un isomorphisme dans $\cA$ pour tout $N \in \cA$.

\clearpage
\printedpage{36}
% head: A. Grothendieck : Dérivateurs, Ch. I, Édition des Universités de Montpellier II et Paris VII
% foot: Version 8 juillet 2001
\noindent\msp{page 62}

\noindent En résumé :

\noindent\textbf{Proposition 1.} \textit{Si $p_*$ et $q_*$ existent, i.e. si $H^*_{\DD}(X, \bullet)$ et $H^*_{\DD}(Y, \bullet)$ existent (il suffit même que ça existe pour des coefficients constants), alors $f$ est une $\DD$-équivalence si et seulement si $f$ induit un isomorphisme pour les $\DD$-cohomologies à coefficients constants de $X$ et $Y$. Si $p_!$ et $q_!$ existent, i.e. si $H^{\DD}_*(X, \bullet)$ et $H^{\DD}_*(Y, \bullet)$ existent, $f$ est une $\DD$-équivalence si et seulement si $f$ induit un isomorphisme pour les $\DD$-homologies à coefficients constants de $X$ et $Y$.}

\noindent\textbf{Remarque 1.} Pour que $f$ soit \texttt{[une]} $\DD$-équivalence, il \textit{suffit} que pour tout $F$ \textit{constant} sur $Y$ le morphisme d'adjonction $F \to f_*f^*F$ soit \texttt{[un]} isomorphisme.

\noindent\textbf{Remarque 2.} Si $f$, $f' : X \to Y$ sont dans une même composante connexe de $\underline{\Hom}(X, Y)$, alors $f$ est $\DD$-équivalence $\Leftrightarrow$ $f'$ $\DD$-équivalence.

\noindent\textbf{Proposition 2.} \textit{Les isomorphismes sont des $\DD$-équivalences. Si $X \overset{f}{\to} Y \overset{g}{\to} Z$, si deux parmis $f$, $g$, $gf$ sont des $\DD$-équivalences, le troisième aussi.} % [sic] parmis -> parmi, spelling

\noindent\textbf{Définition 2.} On dit que $X$ dans $\Dia$ est $\DD$-\textit{asphérique} (ou $\DD$-\textit{contractile}, ou $\DD$-\textit{trivial}) si $X \overset{p}{\to} e$ est une $\DD$-équivalence, \textit{i.e.} si $p^*$ est pleinement fidèle.

Si $p_*$ existe, cela signifie donc que pour $M \in \Ob(\cA)$
\[
M \overset{\mathrm{adj}}{\isolongto} p_*p^*M \quad , \qquad \textit{i.e.} \quad M \isolongto H^*_{\DD}(X, M) \quad ,
\]
et si $p_!$ existe, cela signifie que pour tout $M$
\[
p_!p^*M \overset{\mathrm{adj}}{\isolongto} M \quad , \qquad \textit{i.e.} \quad H^{\DD}_*(X, M) \isolongto M \quad .
\]
\textbf{Corollaire.} \textit{Si $Y$} (resp. $X$) \textit{est contractile, alors $f : X \to Y$ est une $\DD$-équivalence si et seulement si $X$} (resp. $Y$) \textit{est contractile.}

\noindent\textsc{Exemple}. Prenons $\DD = \HOT$. (On peut prendre $\Dia$ = catégorie des petites catégories). Alors, pour $M \in \cA = \Hot$ et $X$ petite catégorie, on a
\[
H^*_{\HOT}(X, M) \simeq \Hom_{\cA}(\hot(X), M) \quad ,
\]
où
\[
\hot : \mathcal{C}\mathit{at} \longrightarrow \cA = \Hot
\]
est le foncteur canonique qui associe à tout $X$ son “type d'homotopie”. On a
\[
H^{\HOT}_*(X, M) = \hot(X) \times M \quad .
\]

\noindent\msp{page 63}

\noindent On en déduit que $f : X \to Y$ est une $\HOT$-équivalence (au sens de la définition 1) si et seulement si c'est une $\hot$-équivalence (ou “équivalence faible”) \textit{i.e.} si et seulement si
\[
\hot(f) : \hot(X) \isolongto \hot(Y)
\]
\printedpage{37}
% head: A. Grothendieck : Dérivateurs, Ch. I, Édition des Universités de Montpellier II et Paris VII
% foot: Version 8 juillet 2001
est un isomorphisme.

\noindent\textbf{Proposition 2.} \textit{Soit $f : X \to Y$. Conditions équivalentes (si $\DD$} \texttt{[est un]} \textit{dérivateur} à gauche). % ?? uncertain: font of "à gauche)." -- upright or italic

a) \textit{$f$ est une $\DD$-équivalence, et le reste après tout changement de base $Y' \to Y$ qui est un isomorphisme local (par exemple pour $Y/F \to Y$, $F \in \Ob \widehat{Y}$).}

b) \textit{Pour tout $y \in \Ob Y$, le foncteur induit $X/y \to Y/y$ est une $\DD$-équivalence.}

c) \textit{Pour tout $y \in \Ob Y$, $X/y$ est $\DD$-contractile.}\footnote[25]{Autres conditions équivalentes :\par
\indent d) \textit{Pour tout faisceau \underline{constant} $M_Y$ sur $Y$, on a $M_Y \isoto f_*f^*M_Y$.}\par
\indent d$'$) \textit{Comme} d) \textit{avec faisceau \underline{localement} constant.}}

\noindent N.B. On dit que $g : Z \to Y$ est un \textit{isomorphisme local} si pour tout $z \in \Ob Z$, posons $y = g(z)$, le morphisme induit sur les catégories localisées en $z$ et $y$
\[
Z/z \isolongto Y/y
\]
est un isomorphisme. Par exemple toute immersion ouverte, tout foncteur induit $Y/y \to Y$ (foncteur de localisation) sont des isomorphismes locaux.

Il en résulte a) $\Rightarrow$ b) trivialement.

b) $\Leftrightarrow$ c). Résulte du corollaire de la page précédente, si l'on sait que $Y/y$ est $\DD$-contractile\footnote[26]{Cf. cor. 1 de \textbf{Der 5} (g)}. Cela résultera de \textbf{Der 5} ci-dessous.

\noindent\msp{page 64}

c) $\Rightarrow$ a) “Passage du local au global”. Comme l'hypothèse c) est stable par changement de base par immersion ouverte, il suffit de prouver que c) implique que $f$ est une $\DD$-équivalence, \textit{i.e.} que pour tout $F = M_Y$ ($M \in \Ob \cA$), coefficient constant sur $Y$, l'homomorphisme d'adjonction
\[
M_Y \longrightarrow f_*f^*M_Y = f_*M_X
\]
est \texttt{[un]} isomorphisme. Il suffit de prouver que c'est un isomorphisme fibre par fibre (\textbf{Der 2}), donc que pour tout $y \in Y$,
\[
M_y \longrightarrow f_*(M_X)_y
\]
est un isomorphisme. Par \textbf{Der 4} (g), on a
\[
f_*(M_X)_y \isolongto \Gamma(X/y, M) \quad ,
\]
et l'hypothèse que $X/y$ est $\DD$-contractile donne ce qu'on veut.

\noindent\textbf{Proposition 2$'$.} \textit{Dualement, si $\DD$ est un dérivateur à gauche} \texttt{[plutôt à droite]}\textit{, conditions équivalentes sur $f : X \to Y$ :}

a) \textit{$f$ $\DD$-équivalence,} \texttt{[et]} \textit{le reste par tout changement de base $g : Y' \to Y$ qui est un “isomorphisme colocal” (i.e. pour tout $y' \in \Ob Y'$, le foncteur induit sur les catégories “colocalisées” $y'\backslash Y' \isolongto y\backslash Y$ (où $y = g(y')$)} \texttt{[est un isomorphisme]}\textit{).}

b) \textit{Pour tout $y \in Y$, le foncteur induit $y\backslash X \to y\backslash Y$ est une $\DD$-équivalence.}

c) \textit{Pour tout $y \in Y$, $y\backslash Y$ est $\DD$-contractile.}

\noindent Pour le prouver, on doit utiliser que les $y\backslash Y$ sont $\DD$-contractiles (cf. cor. 1 de \textbf{Der 5} (d)).

\clearpage
\printedpage{38}
% head: A. Grothendieck : Dérivateurs, Ch. I, Édition des Universités de Montpellier II et Paris VII
% foot: Version 8 juillet 2001
\noindent\msp{page 65}

\noindent\textbf{Définition 3.} Soit $\DD$ un prédérivateur (peut-être même un dérivateur à gauche ou à droite). On dit que $f$ est une $\DD$-équivalence \textit{stricte} si $f$ satisfait la condition a) de \texttt{[la]} proposition 2. On dit que c'est une $\DD$-équivalence \textit{costricte} si $f$ satisfait celle de \texttt{[la]} proposition $2'$; on dit aussi parfois que $f$ est $\DD$-cofinal. (La condition duale, \textit{i.e.} $\DD$-équivalence stricte, s'appelerait co-$\DD$-cofinal!)\footnote[27]{Meilleure terminologie : $f$ est $\DD$-\textit{asphérique} ($\DD$-équivalence stricte), $\DD$-\textit{coasphérique} ($\DD$-équivalence co-stricte).} % [sic] s'appelerait -> s'appellerait, spelling

\noindent\textbf{Proposition 3.} \textit{Soit $X \overset{f}{\to} Y \overset{g}{\to} Z$. Si $f, g$ sont des $\DD$-équivalences strictes, alors $gf$ aussi. Si $gf$ et $f$ sont des $\DD$-équivalences strictes, alors $g$ aussi. Si $\DD$} \texttt{[est un]} \textit{dérivateur à gauche ou à droite, toute équivalence} \texttt{[de catégories]} \textit{est une $\DD$-équivalence stricte.} (Car elle a un adjoint à droite.)

On utilise le fait que la notion d'isomorphisme local est stable par changement de base.

\noindent\msp{page 66}

\noindent\textbf{Proposition 4.} \textit{Soient $\DD$ un dérivateur à droite, $f : X \to Y$ un foncteur dans} $\Dia$\textit{, tel que pour $y \in Y$, $y\backslash X$ soit $\DD$-contractile (donc $f$ une $\DD$-équivalence costricte, sous réserve que $\DD$ soit un dérivateur à droite \ldots\ ). Alors pour tout $F$ dans $\DD(Y)$ (pas nécessairement à coefficients constant) on a} % [sic] constant -> constants, agreement
\[
H^*_{\DD}(Y, F) \isolongto H^*_{\DD}(X, f^*(F))
\]
\textit{est un isomorphisme. Réciproquement, si la condition est vraie pour tout $F$, alors les $y\backslash X$ sont $\DD$-contractiles.}

\noindent\textsc{Démonstration}. Il revient au même de dire que pour tout $M \in \Ob \cA = \DD(e)$, on a
\[
\Hom(M, H^*_{\DD}(Y, F)) \overset{\mathrm{iso}}{\longrightarrow} \Hom(M, H^*_{\DD}(X, f^*(F))) \ .
\]
Mais par adjonction, cette flèche s'interprète comme
\[
\Hom(M_Y, F) \longrightarrow \Hom(M_X, f^*F) \tag*{$(*)$}
\]
D'autre part, si $f_!$ existe, on a
\[
\Hom(M_X, f^*F) = \Hom(f_!(M_X), F) \ ,
\]
et l'homomorphisme $(*)$ est déduit de l'homomorphisme d'adjonction
\[
f_!(M_X) = f_!(f^*(M_Y)) \longrightarrow M_Y \ , \tag*{$(**)$}
\]
en prenant les $\Hom$ dans $F$. Donc la condition équivaut à dire que cet homomorphisme est \texttt{[un]} isomorphisme, quel que soit le coefficient constant $M_Y$ sur $Y$. Pour voir si c'est un isomorphisme, on regarde fibre par fibre (\textbf{Der 2}), or les fibres s'interprètent comme $H^{\DD}_*(y\backslash X, M_X)$ (par \textbf{Der 4} (d)) et comme $M$.

\end{document}
